\section{Component Diagram}

Rispetto alla precedente iterazione, l'architettura del sistema è stata estesa 
mantenendo invariata la struttura a più livelli già adottata. La separazione tra 
interfaccia utente, logica applicativa e gestione dei dati è stata preservata, 
introducendo nuove funzionalità sia lato front-end che lato back-end.

\begin{figure}[H]
\centering
\includegraphics[width=\linewidth]{Iterazione3/Diagrammi/ComponentDiagramIt3.png}
\caption{Diagramma dei componenti}
\label{fig:ComponentDiagramIterazione3}
\end{figure}

In questa iterazione è stata aggiunta una nuova pagina nel livello \textit{View}, 
ovvero \textit{ProfilePage}, che consente all'utente di visualizzare e modificare 
il proprio profilo, aggiornare le credenziali e eliminare il proprio account. 
Questa pagina interagisce con due controller già esistenti: \textit{UserController}, 
tramite l'interfaccia \texttt{api/user}, per le operazioni sul profilo, e 
\textit{AuthController}, tramite \texttt{api/auth}, per la modifica della password.

Nel livello \textit{Service} è stato introdotto il nuovo componente 
\textit{ExportService}, esposto attraverso l'interfaccia \textit{ExportServiceIF}. 
Questo service è responsabile della generazione del PDF riepilogativo di un 
itinerario ed è invocato direttamente da \textit{TripController}. 


Il diagramma mostra quindi un'evoluzione coerente dell'architettura: la nuova 
pagina del front-end sfrutta i controller esistenti, il nuovo service di esportazione 
si integra nel layer applicativo senza alterare le responsabilità dei componenti 
preesistenti, e le dipendenze aggiornate riflettono i nuovi requisiti funzionali 
introdotti in questa iterazione.